% Chen Yang's CV, adapted from Sina Atalay's RenderCV EngineeringResumes theme:
% https://www.overleaf.com/latex/templates/rendercv-engineeringresumes-theme/shwqvsxdgkjy
% Keep this source and public/cv.pdf in sync with the build command in cv/README.md.
\documentclass[10pt,letterpaper]{article}

\usepackage[ignoreheadfoot,top=1.4cm,bottom=1.35cm,left=1.85cm,right=1.85cm]{geometry}
\usepackage{titlesec}
\usepackage{tabularx}
\usepackage{array}
\usepackage{xcolor}
\usepackage{enumitem}
\usepackage{charter}
\usepackage{changepage}
\usepackage{paracol}
\usepackage{needspace}
\usepackage[T1]{fontenc}
\usepackage[utf8]{inputenc}
\input{glyphtounicode}
\pdfgentounicode=1
\definecolor{primaryColor}{RGB}{0,0,0}
\usepackage[pdftitle={Chen Yang - Academic CV},pdfauthor={Chen Yang},pdfcreator={LaTeX with RenderCV EngineeringResumes theme},colorlinks=true,urlcolor=primaryColor]{hyperref}

\raggedright
\linespread{1.04}
\pagestyle{empty}
\setcounter{secnumdepth}{0}
\setlength{\parindent}{0pt}
\setlength{\topskip}{0pt}
\setlength{\columnsep}{0.4cm}
\titleformat{\section}{\needspace{4\baselineskip}\bfseries\large}{}{0pt}{}[\vspace{1pt}\titlerule]
\titlespacing{\section}{-1pt}{0.23cm}{0.18cm}
\renewcommand\labelitemi{$\vcenter{\hbox{\small$\bullet$}}$}

\newenvironment{highlights}{\begin{itemize}[topsep=0.09cm,parsep=0pt,partopsep=0pt,itemsep=0.04cm,leftmargin=0.4cm]}{\end{itemize}}
\newenvironment{onecolentry}{\begin{adjustwidth}{0.00001cm}{0.00001cm}}{\end{adjustwidth}}
\newenvironment{twocolentry}[2][]{%
  \onecolentry
  \def\secondColumn{#2}
  \setcolumnwidth{\fill,3.7cm}
  \begin{paracol}{2}
}{%
  \switchcolumn\raggedleft\secondColumn
  \end{paracol}
  \endonecolentry
}
\newenvironment{header}{\setlength{\topsep}{0pt}\par\kern\topsep\centering\linespread{1.5}}{\par\kern\topsep}

\begin{document}
\begin{header}
  \fontsize{25pt}{25pt}\selectfont Chen Yang

  \vspace{5pt}
  \normalsize
  \href{mailto:chenyang@link.cuhk.edu.cn}{chenyang@link.cuhk.edu.cn}
  \quad|\quad
  (+86) 15267711552\\
  The Chinese University of Hong Kong, Shenzhen, China, 518172
\end{header}

\vspace{0.08cm}

\section{Education}
\begin{twocolentry}{09/2022 -- 05/2027 (Expected)}
  \textbf{The Chinese University of Hong Kong, Shenzhen}\\
  Major in Electrical and Computer Engineering
\end{twocolentry}
\begin{onecolentry}
  \begin{highlights}
    \item Cumulative GPA: 3.856/4.0 (Rank: 5/116)
    \item Research Interests: Robot Learning, Reinforcement Learning
    \item Awards \& Honors: Creativity and Innovation Award, 2024; Academic Scholarship, 2023 \& 2024;\\
      Dean's List, 2023 \& 2024 \& 2025 \& 2026
  \end{highlights}
\end{onecolentry}
\vspace{0.12cm}
\begin{twocolentry}{09/2024 -- 12/2024}
  \textbf{University of California, Berkeley}\\
  Visiting Student
\end{twocolentry}
\begin{onecolentry}
  \begin{highlights}
    \item Cumulative GPA: 4.0/4.0
    \item Related Courses: Computer Architecture, Artificial Intelligence, Discrete Math (A+)
  \end{highlights}
\end{onecolentry}

\section{Publications}
\begin{onecolentry}
  \begin{enumerate}[label={[\arabic*]},leftmargin=0.55cm,topsep=0.03cm,parsep=0pt,itemsep=0.04cm]
    \item Feiyang Wu, Chenxiao Gao, \textbf{Chen Yang}, Ye Zhao, Bo Dai, and Anqi Wu. ``\href{https://arxiv.org/abs/2609.37677}{Learning Expressive and Compositional Motion Representation via Spectral Skills}.'' 2026. Submitted to ICLR; IROS 2026 Workshop Best Paper.
  \end{enumerate}
\end{onecolentry}

\section{Research Experience}
\begin{twocolentry}{01/2026 -- Present}
  \textbf{Robot Motion Generation}\\
  Research Intern, Advisor: Guiliang Liu\\
  DexForce
\end{twocolentry}
\begin{onecolentry}
  \begin{highlights}
    \item Lead research on robot motion generation, focusing on producing feasible motions under diverse constraints.
    \item Further details coming soon.
  \end{highlights}
\end{onecolentry}
\vspace{0.16cm}

\begin{twocolentry}{06/2026 -- 09/2026}
  \textbf{Learning Expressive and Compositional Motion Representation via Spectral Skills}\\
  Research Intern, Advisors: Prof. Ye Zhao and Prof. Anqi Wu\\
  LiDAR Lab \& BRAINML Lab, Georgia Institute of Technology
\end{twocolentry}
\begin{onecolentry}
  \begin{highlights}
    \item Pretrained, fine-tuned, and evaluated the high-level planner using predictive spectral skills as commands for a low-level whole-body controller.
    \item Reproduced baselines, including SONIC, and helped design ablation studies to assess key modules and design choices.
    \item Worked with the team to build sim-to-sim and sim-to-real pipelines and deploy simulation-trained policies on a Unitree G1 humanoid.
  \end{highlights}
\end{onecolentry}
\vspace{0.16cm}

\begin{twocolentry}{09/2023 -- 04/2024}
  \textbf{Smart Stop Snoring Pillow}\\
  Research Assistant, Advisor: Prof. Jian Zhu\\
  Soft Robotics Lab, CUHKSZ
\end{twocolentry}
\begin{onecolentry}
  \begin{highlights}
    \item Designed and built an integrated pneumatic soft-robotic pillow for non-invasive snoring mitigation, combining pneumatic, electronic, sensing, and embedded-control subsystems.
    \item Modeled airflow through the tubing and airbag using Poiseuille’s law and pressure differentials from two barometers to estimate flow and precisely control airbag height; built ROS-based middleware connecting the host computer and microcontroller.
    \item Applied FFT to barometric-pressure signals to estimate heart and respiratory rates, and developed a Transformer-CNN classifier to distinguish snoring from non-snoring events.
  \end{highlights}
\end{onecolentry}
\vspace{0.16cm}

\newpage
\section{Internship}
\begin{twocolentry}{01/2026 -- Present}
  \textbf{DexForce}\\
  Research Intern
\end{twocolentry}
\begin{onecolentry}
  \begin{highlights}
    \item Core contributor to \href{https://github.com/DexForce/EmbodiChain}{EmbodiChain}, an end-to-end, GPU-accelerated, modular platform for embodied AI.
    \item Refined and developed EmbodiChain's RL framework for manipulation, humanoid motion, and VLA post-training.
  \end{highlights}
\end{onecolentry}
\vspace{0.16cm}

\begin{twocolentry}{04/2024 -- 08/2024}
  \textbf{Shenzhen Research Institute of Big Data}\\
  Research Intern
\end{twocolentry}
\begin{onecolentry}
  \begin{highlights}
    \item Developed an LSTM-T-KAN model for efficient, accurate long-horizon building-load forecasting to support power allocation.
    \item Extracted frequency-domain features using the Fast Fourier Transform (FFT).
  \end{highlights}
\end{onecolentry}

\section{Teaching}
\begin{onecolentry}
  \textbf{Teaching Assistant}
\end{onecolentry}
\vspace{0.08cm}
\begin{twocolentry}{2026--27 Term 1}
  ECE3080 Introduction to Embedded Systems
\end{twocolentry}
\vspace{0.08cm}
\begin{twocolentry}{2026--27 Term 1}
  ECE3810 Embedded Systems Laboratory
\end{twocolentry}
\vspace{0.08cm}
\begin{twocolentry}{2023--24 Term 2}
  PHY1001 Mechanics
\end{twocolentry}

\section{Skills}
\begin{onecolentry}
  \textbf{Programming Languages:} Python, C/C++, MATLAB, Java\\
  \textbf{Technologies \& Frameworks:} PyTorch, TensorFlow, ROS, Linux, Git\\
  \textbf{Languages:} English (TOEFL: 102), Chinese (Native)
\end{onecolentry}

\end{document}
